\documentclass[spanish,letterpaper,oneside]{article}

\def\documenttitle {Foro de semana 1: Antecedentes de los derechos
humanos}
\def\documentsubtitle {Participacion del foro diagnostico de semana 1}
\def\documentsubject {Semana 1 - Antecedentes de los derechos humanos}
\def\documentauthor {Martin Jonathan de la Cruz}
\def\coursename {Antecedentes de los derechos humanos}
\def\coursecode {005}
\def\universityname {Universidad Abierta y a Distancia de Mexico}
\def\universityfaculty {Licenciatura en Derecho}
\def\universitydepartment {Antecedentes de los derechos humanos}
\def\universitydepartmentimage {departamentos/UnADM}
\def\universitydepartmentimagecfg {height=1.57cm}
\def\universitylocation {Roma Norte, Ciudad de Mexico}
\def\authortable {
  \begin{tabular}{ll}
    \textbf{Alumno:} & Martin Jonathan de la Cruz \\
    Matricula: & ES2611202040 \\
    Figura docente: & Daniel Alexis Lozano Ortega \\
    Semestre/Bloque: & 2 / 2 \\
    Estado: & Borrador para revision \\
    \multicolumn{2}{l}{Fecha de realizacion: \today} \\
    \multicolumn{2}{l}{\universitylocation}
  \end{tabular}
}

\input{template}
\usepackage{helvet}
\renewcommand{\familydefault}{\sfdefault}
\usepackage{array}
\usepackage{booktabs}
\usepackage{longtable}
\usepackage{calc}
\usepackage{tikz}
\usetikzlibrary{arrows.meta,positioning,calc}
\setcitestyle{authoryear,open={(},close={)}}
\providecommand{\tightlist}{\setlength{\itemsep}{0pt}\setlength{\parskip}{0pt}}
\providecommand{\pandocbounded}[1]{#1}
\providecommand{\passthrough}[1]{#1}
\setlength{\emergencystretch}{3em}

\usepackage[most]{tcolorbox}
\definecolor{unadmGreenDark}{HTML}{174A3A}
\definecolor{unadmGreen}{HTML}{5F8F3A}
\definecolor{unadmPaper}{HTML}{F6F7F2}
\newtcolorbox{forobox}[1][]{%
  enhanced, breakable,
  colback=unadmPaper, colframe=unadmGreen,
  boxrule=0pt, leftrule=3pt, arc=1.2mm,
  left=4mm, right=4mm, top=3mm, bottom=3mm,
  fonttitle=\bfseries\color{white}, coltitle=white, colbacktitle=unadmGreenDark,
  attach boxed title to top left={xshift=4mm,yshift=-2mm},
  boxed title style={colframe=unadmGreenDark,arc=0.6mm},
  #1
}

\def\coverwatermarkenabled {true}
\def\coverwatermarkimage {img/departamentos/UnADM.pdf}
\def\coverwatermarkopacity {0.16}
\def\coverwatermarkwidth {0.72\paperwidth}
\def\coverwatermarkxshift {0cm}
\def\coverwatermarkyshift {-0.35cm}
\newcommand{\insertcoverwatermark}{%
  \ifthenelse{\equal{\coverwatermarkenabled}{true}}{%
    \AddToShipoutPictureBG*{%
      \begin{tikzpicture}[remember picture,overlay]
        \node[opacity=\coverwatermarkopacity,inner sep=0pt,
          xshift=\coverwatermarkxshift,yshift=\coverwatermarkyshift]
          at (current page.center) {%
            \includegraphics[width=\coverwatermarkwidth]{\coverwatermarkimage}%
          };
      \end{tikzpicture}%
    }%
  }{}%
}

\begin{document}
\insertcoverwatermark
\templatePortrait
\templatePagecfg
\onehalfspacing
\templateIndex
\templateFinalcfg

\section{Identificacion y estado}

Propuesta de participacion basada en el proyecto, elaborada con apoyo de GitHub Copilot. Requiere revision personal; no publicada. No acredita conocimientos previos ni experiencias personales confirmadas. Edad y residencia, cuando se solicitan, permanecen pendientes.

\section{Participacion principal}
\begin{forobox}[title={Participacion principal -- propuesta no publicada}]

\textbf{¿Que entiendes por derechos humanos?}

Entiendo que son derechos que corresponden a todas las personas por su
dignidad y que deben respetarse sin discriminacion. Su importancia esta
en proteger la libertad, la igualdad y las condiciones necesarias para
vivir y desarrollarse. Considero que su reconocimiento debe acompanarse
de mecanismos que permitan exigir su respeto cuando una autoridad o una
persona los vulnera.

\textbf{Menciona al menos tres derechos humanos que conozcas.}

El derecho a la educacion, el derecho a la salud y el derecho a la
igualdad y no discriminacion.

\textbf{¿Que tema relacionado con los derechos humanos te gustaria
aprender durante esta asignatura?}

Me interesa comprender como se paso del reconocimiento de ciertos
derechos para grupos limitados a la idea de derechos para todas las
personas. Tambien quiero conocer como esa evolucion se relaciona con las
garantias constitucionales en Mexico y con los mecanismos que permiten
proteger los derechos en la practica.

\end{forobox}

\section{Respuestas y alcance}

No se agregan respuestas ficticias. La consigna consultada no exige replicas; Antecedentes indica expresamente que no es necesario responder a companeros. Si el docente solicita respuestas adicionales, deben documentarse en cajas independientes tras leer las intervenciones correspondientes.

\end{document}